\begin{otherlanguage}{english}
\chapter*{\centering Abstract}
Retail networks must select replenishment policies for heterogeneous demand. This dissertation proposes the Adaptive Inventory Policy Engine (AIPE), integrating data preparation, characterization, calibration, and contextual selection. Its Policy Selection Engine (PSE) either maps a given profile to a policy or predicts the policy for each series using supervised learning. Seven candidates are fitted on training data, selected using validation data from other stores, and evaluated on held-out future periods. In Bahia, 145 series at two origins yield 4,350 fits. Profile-based selection increases mean cost by 1.93\% and mean fill rate by 3.89 percentage points, reducing cases below 0.70 from 119 to 100 among 283 positive-demand cases. In the M5 pilot, 30 series from ten stores at three origins yield 1,350 fits. Predictive selection costs 13.9\% more and serves 4.11 percentage points less than global selection; explicitly adding the profile does not change its recommendations. Limited class support and unstable labels constrain learning. An exploratory fallback to profile-based selection reduces the predictive selector's loss without demonstrating superiority over the references. Retrospective portfolio oracles and perfect-information planning separate selection errors, portfolio gaps, and initial physical limitations. The contribution is a framework and a traceable evaluation that reports both benefits and limitations. Retrospective cohorts, dataset-specific calibrations, short horizons, and development on the pilot require independent confirmation before claiming a general AIPE advantage.

\vspace{2em}
\textbf{Keywords}: Inventory management; Policy selection; Intermittent demand; Supervised learning; Service level.
\end{otherlanguage}
